%
% display_lifecycle.tex
%
% A Z specification of the lux display singleton lifecycle: concurrent
% agents contending for one Unix-socket path under a two-lock discipline
% (spawn lock + bind lock).  Formalises the design hardened empirically
% across lux-w8t5 and lux-h29e.
%
% Source of record:
%   src/punt_lux/paths.py         -- DisplayPaths: probe, spawn/bind locks,
%                                    cleanup_stale, reap
%   src/punt_lux/socket_server.py -- SocketServer.setup: the
%                                    probe -> cleanup -> bind -> listen
%                                    critical section under bind_lock
%
% Type-check:  fuzz -t docs/display_lifecycle.tex
% Model-check (see Section "Verification" for the exact goal predicates):
%   probcli docs/display_lifecycle.tex -model_check -p DEFAULT_SETSIZE 2
%   probcli docs/display_lifecycle.tex -model_check -nodead \
%       -goal "card({a|a:AGENT & phase(a):{bound,listening,serving}})>1" \
%       -p DEFAULT_SETSIZE 2
%
% Formatting follows the fuzz house rules: \quad~ for continuation
% indent (never \t1..\t9), schema lines under 80 columns, predicates
% broken at \land/\lor/\implies with the operator starting the
% continuation line.
%

\documentclass[a4paper,10pt,fleqn]{article}
\usepackage[margin=1in]{geometry}
\usepackage{fuzz}
\usepackage[colorlinks=true,linkcolor=blue,citecolor=blue,urlcolor=blue]{hyperref}

\begin{document}

\title{The lux Display Singleton Lifecycle: a Z Specification}
\author{Formal model of the two-lock bind-race discipline}
\date{July 2026}
\maketitle

% ============================================================
\section{Introduction}
% ============================================================

Two or more agents may, at the same instant, decide that no display is
running and try to become the one display that owns a single Unix-socket
path.  Exactly one must win.  The mechanism that arbitrates the race is a
pair of advisory file locks --- a \emph{spawn lock} held end-to-end by
\texttt{ensure} and \texttt{reap}, and a \emph{bind lock} held by
\texttt{setup} across the whole
\texttt{probe $\to$ cleanup $\to$ bind $\to$ listen} critical section.

The subtlety the design must survive is the \emph{bind window}.  A socket
that has been \texttt{bind}-ed but not yet \texttt{listen}-ed refuses
connections; a liveness probe therefore reads it as \emph{dead}.  Without
serialisation, a second agent probing during that window concludes the
path is stale, unlinks the first agent's fresh socket, and binds its own
--- and now two displays believe they own the path.  That is the round-2
defect fixed on branch \texttt{fix/lux-h29e-bind-race}.

This document models the lifecycle as a finite state machine over a small
fixed set of agents, so that ProB can enumerate every interleaving and
either confirm the five invariants or produce the exact offending trace.

% ============================================================
\section{Basic Types}
% ============================================================

\begin{zed}
  [AGENT]
\end{zed}

\texttt{AGENT} is the set of concurrent processes contending for the
path.  Two agents suffice to exhibit the two-winner race; three exercise
lock contention more thoroughly.  The carrier is bounded by ProB's
\texttt{DEFAULT\_SETSIZE}, so no explicit cardinality constant is needed.

% ============================================================
\section{Free Types}
% ============================================================

\begin{zed}
  SOCKSTATUS ::= sabsent | sbound | slistening | sstale
\end{zed}

The status of the \emph{file at the socket path}, as a liveness probe
would observe it.  \texttt{sabsent}: no file.  \texttt{sbound}: a socket
is bound but \texttt{listen} has not been called --- the bind window; a
probe reads it dead.  \texttt{slistening}: a live owner is accepting
connections.  \texttt{sstale}: a leftover socket file whose owner has
died --- a probe reads it dead.

\begin{zed}
  PHASE ::= \\
  \quad~ idle | hasSpawn | hasBind | probed | bound | \\
  \quad~ listening | serving | lostRace | exited | reaping
\end{zed}

Each agent's own progress.  \texttt{idle} before it acts;
\texttt{hasSpawn}/\texttt{hasBind} once it holds the corresponding lock;
\texttt{probed} after it has observed the path; \texttt{bound} after a
successful \texttt{bind}; \texttt{listening} after \texttt{listen};
\texttt{serving} once it has released the bind lock and is the running
display; \texttt{lostRace} for an agent that observed a live owner or
lost the bind; \texttt{exited} for a terminated agent; \texttt{reaping}
for a reaper mid-teardown.

% ============================================================
\section{State}
% ============================================================

The state is flat: one schema, five components.  Its predicate carries
only \emph{structural} invariants --- cardinality bounds on the locks and
owner, and the correspondence between the path's status and whether it
has an owner.  These hold in every design, correct or buggy, so they act
as coherence constraints, not as the safety properties under test.  The
safety properties (Section~\ref{sec:inv}) are deliberately \emph{absent}
here, so that a violating state remains reachable and ProB can find it.

\begin{schema}{Sys}
  socketStatus : SOCKSTATUS \\
  owner : \power AGENT \\
  spawnLock : \power AGENT \\
  bindLock : \power AGENT \\
  phase : AGENT \fun PHASE
\where
  \# owner \leq 1 \\
  \# spawnLock \leq 1 \\
  \# bindLock \leq 1 \\
  socketStatus = sabsent \implies owner = \emptyset \\
  socketStatus = sstale \implies owner = \emptyset \\
  socketStatus = sbound \implies \# owner = 1 \\
  socketStatus = slistening \implies \# owner = 1
\end{schema}

\texttt{owner} is the singleton (or empty) set naming the agent whose
socket is currently installed at the path; a lock component names its
current holder.  Modelling these as sets of size at most one lets the
empty set stand for ``none'' without an extra sentinel value.

% ============================================================
\section{Initialization}
% ============================================================

\begin{schema}{Init}
  Sys'
\where
  socketStatus' = sabsent \\
  owner' = \emptyset \\
  spawnLock' = \emptyset \\
  bindLock' = \emptyset \\
  phase' = (\lambda a : AGENT @ idle)
\end{schema}

A single initial state: no file, no locks held, every agent idle.

% ============================================================
\section{Lock Operations}
% ============================================================

The two locks are acquired non-blockingly: an acquire is \emph{enabled}
only when the lock is free, so a ``blocked'' agent is simply an agent
whose acquire is not enabled.  Cyclic waiting --- and hence deadlock ---
is what ProB's deadlock check must rule out.

The lock order is spawn-outer, bind-inner.  An agent that already holds
the bind lock never asks for the spawn lock, and vice versa: each agent
holds at most one of the two.  This is encoded directly in the acquire
guards, and is the syntactic reason no cyclic wait can form.

\begin{schema}{AcquireBindLock}
  \Delta Sys \\
  a? : AGENT
\where
  phase~a? = idle \\
  a? \notin spawnLock \\
  bindLock = \emptyset \\
  bindLock' = \{ a? \} \\
  phase' = phase \oplus \{ a? \mapsto hasBind \} \\
  socketStatus' = socketStatus \\
  owner' = owner \\
  spawnLock' = spawnLock
\end{schema}

\begin{schema}{ReleaseBindLock}
  \Delta Sys \\
  a? : AGENT
\where
  a? \in bindLock \\
  phase~a? = listening \\
  bindLock' = \emptyset \\
  phase' = phase \oplus \{ a? \mapsto serving \} \\
  socketStatus' = socketStatus \\
  owner' = owner \\
  spawnLock' = spawnLock
\end{schema}

Releasing the bind lock is the last step of the critical section: the
agent has bound and listened, and now becomes the running display.

\begin{schema}{AcquireSpawnLock}
  \Delta Sys \\
  a? : AGENT
\where
  phase~a? = idle \\
  a? \notin bindLock \\
  spawnLock = \emptyset \\
  spawnLock' = \{ a? \} \\
  phase' = phase \oplus \{ a? \mapsto hasSpawn \} \\
  socketStatus' = socketStatus \\
  owner' = owner \\
  bindLock' = bindLock
\end{schema}

\begin{schema}{ReleaseSpawnLock}
  \Delta Sys \\
  a? : AGENT
\where
  a? \in spawnLock \\
  phase~a? = reaping \\
  spawnLock' = \emptyset \\
  phase' = phase \oplus \{ a? \mapsto exited \} \\
  socketStatus' = socketStatus \\
  owner' = owner \\
  bindLock' = bindLock
\end{schema}

% ============================================================
\section{The Bind Critical Section}
% ============================================================

Every operation in this section is guarded by \texttt{a? $\in$ bindLock}:
in the current design the probe, the stale cleanup, the bind, and the
listen all run under the one held bind lock.  That single guard is the
whole serialisation; removing it from \texttt{CleanupStale} and
\texttt{Bind} yields the buggy variant of Section~\ref{sec:fidelity}.

\texttt{Probe} records that the agent has observed the path.  Observation
and the action taken on it are separate steps, so the state may change
between them --- this is the time-of-check/time-of-use gap the bind lock
must close.

\begin{schema}{Probe}
  \Delta Sys \\
  a? : AGENT
\where
  a? \in bindLock \\
  phase~a? = hasBind \\
  phase' = phase \oplus \{ a? \mapsto probed \} \\
  socketStatus' = socketStatus \\
  owner' = owner \\
  spawnLock' = spawnLock \\
  bindLock' = bindLock
\end{schema}

\texttt{CleanupStale} unlinks a socket that probes \emph{dead}: a stale
leftover (\texttt{sstale}) or --- and this is the dangerous case --- one
that is bound but not yet listening (\texttt{sbound}).  It never touches
a live (\texttt{slistening}) socket.  A probing agent never owns the path
(it has not bound), so this only ever removes another agent's file.

\begin{schema}{CleanupStale}
  \Delta Sys \\
  a? : AGENT
\where
  a? \in bindLock \\
  phase~a? = probed \\
  socketStatus \in \{ sstale, sbound \} \\
  socketStatus' = sabsent \\
  owner' = \emptyset \\
  phase' = phase \\
  spawnLock' = spawnLock \\
  bindLock' = bindLock
\end{schema}

Under the bind lock this operation's \texttt{sbound} branch is
unreachable: the only agent that could hold the lock is \texttt{a?}
itself, and it has not bound, so no socket is in the bind window.  ProB
confirms the branch never fires (Section~\ref{sec:verify}).  In the
buggy variant it fires, and that is precisely the two-winner defect.

\begin{schema}{Bind}
  \Delta Sys \\
  a? : AGENT
\where
  a? \in bindLock \\
  phase~a? = probed \\
  socketStatus = sabsent \\
  socketStatus' = sbound \\
  owner' = \{ a? \} \\
  phase' = phase \oplus \{ a? \mapsto bound \} \\
  spawnLock' = spawnLock \\
  bindLock' = bindLock
\end{schema}

\texttt{Bind} is atomic and succeeds only from \texttt{sabsent}.  A
concurrent binder that finds the path already taken loses the race:

\begin{schema}{BindFail}
  \Delta Sys \\
  a? : AGENT
\where
  a? \in bindLock \\
  phase~a? = probed \\
  socketStatus = sbound \\
  a? \notin owner \\
  phase' = phase \oplus \{ a? \mapsto lostRace \} \\
  bindLock' = \emptyset \\
  socketStatus' = socketStatus \\
  owner' = owner \\
  spawnLock' = spawnLock
\end{schema}

\texttt{Listen} promotes the bound socket to live.  An agent listens on
its \emph{own} descriptor: if it still owns the path the path goes live;
if its file was unlinked from under it (the buggy variant), it listens on
an orphaned descriptor and the path's status is unchanged.

\begin{schema}{Listen}
  \Delta Sys \\
  a? : AGENT
\where
  a? \in bindLock \\
  phase~a? = bound \\
  phase' = phase \oplus \{ a? \mapsto listening \} \\
  socketStatus' = \\
  \quad~ \IF a? \in owner \THEN slistening \ELSE socketStatus \\
  owner' = owner \\
  spawnLock' = spawnLock \\
  bindLock' = bindLock
\end{schema}

An agent that probed a \emph{live} path concedes without binding:

\begin{schema}{LoseRaceLive}
  \Delta Sys \\
  a? : AGENT
\where
  a? \in bindLock \\
  phase~a? = probed \\
  socketStatus = slistening \\
  phase' = phase \oplus \{ a? \mapsto lostRace \} \\
  bindLock' = \emptyset \\
  socketStatus' = socketStatus \\
  owner' = owner \\
  spawnLock' = spawnLock
\end{schema}

A lost racer exits, changing nothing but its own phase:

\begin{schema}{Exit}
  \Delta Sys \\
  a? : AGENT
\where
  phase~a? = lostRace \\
  phase' = phase \oplus \{ a? \mapsto exited \} \\
  socketStatus' = socketStatus \\
  owner' = owner \\
  spawnLock' = spawnLock \\
  bindLock' = bindLock
\end{schema}

% ============================================================
\section{Failure and Teardown}
% ============================================================

An owner may crash, leaving a stale file.  Killing the process releases
any lock it held:

\begin{schema}{Crash}
  \Delta Sys \\
  a? : AGENT
\where
  a? \in owner \\
  phase~a? \in \{ listening, serving \} \\
  socketStatus' = sstale \\
  owner' = \emptyset \\
  phase' = phase \oplus \{ a? \mapsto exited \} \\
  bindLock' = bindLock \setminus \{ a? \} \\
  spawnLock' = spawnLock
\end{schema}

A reaper terminates the live owner under the spawn lock, then clears the
files.  The file-clearing step needs the bind lock free (it mirrors
\texttt{reap}'s \texttt{\_clear\_dead\_files\_locked}), so a concurrent
binder is never unlinked mid-bind:

\begin{schema}{Reap}
  \Delta Sys \\
  a? : AGENT
\where
  a? \in spawnLock \\
  phase~a? = hasSpawn \\
  socketStatus = slistening \\
  bindLock = \emptyset \\
  phase' = \\
  \quad~ (phase \oplus (owner \cross \{ exited \})) \\
  \quad~ \oplus \{ a? \mapsto reaping \} \\
  socketStatus' = sabsent \\
  owner' = \emptyset \\
  spawnLock' = spawnLock \\
  bindLock' = bindLock
\end{schema}

When there is no live owner, the reaper simply clears any dead file,
again only with the bind lock free:

\begin{schema}{ReapDead}
  \Delta Sys \\
  a? : AGENT
\where
  a? \in spawnLock \\
  phase~a? = hasSpawn \\
  socketStatus \neq slistening \\
  bindLock = \emptyset \\
  socketStatus' = sabsent \\
  owner' = \emptyset \\
  phase' = phase \oplus \{ a? \mapsto reaping \} \\
  spawnLock' = spawnLock \\
  bindLock' = bindLock
\end{schema}

An exited agent may re-enter as a fresh process.  This keeps the system
live --- it never runs out of enabled operations at a benign terminal
state --- so that the deadlock check reports only genuine lock
deadlocks, never mere termination:

\begin{schema}{Reset}
  \Delta Sys \\
  a? : AGENT
\where
  phase~a? = exited \\
  a? \notin spawnLock \\
  a? \notin bindLock \\
  phase' = phase \oplus \{ a? \mapsto idle \} \\
  socketStatus' = socketStatus \\
  owner' = owner \\
  spawnLock' = spawnLock \\
  bindLock' = bindLock
\end{schema}

% ============================================================
\section{Invariants}
\label{sec:inv}
% ============================================================

Five properties must hold across all reachable states.  They are stated
here as predicates over \texttt{Sys}; how each is discharged is given in
Section~\ref{sec:verify}.  None is placed in the \texttt{Sys} schema: a
predicate placed there would be enforced as a guard on every successor,
silently disabling any operation that would break it and thereby masking
the very defect this model exists to catch.

\paragraph{Invariant 1 --- singleton serving.}
At most one agent is in \texttt{serving}:
\[
  \# \{ a : AGENT | phase~a = serving \} \leq 1.
\]

\paragraph{Invariant 2 --- never unlink live.}
No operation carries \texttt{socketStatus} from \texttt{slistening} to
\texttt{sabsent} except a legitimate \texttt{Reap} of that owner.  This is
a property of the transition relation, established by construction: the
only operations that assign \texttt{socketStatus' = sabsent} are
\texttt{CleanupStale} (guarded \texttt{socketStatus $\in$ \{sstale,
sbound\}}), \texttt{Bind}'s precondition path (from \texttt{sabsent}),
\texttt{Reap} (guarded \texttt{socketStatus = slistening}), and
\texttt{ReapDead} (guarded \texttt{socketStatus $\neq$ slistening}).  Of
these, only \texttt{Reap} starts from \texttt{slistening}, and it does so
only after terminating the owner.  \texttt{CleanupStale}'s guard
excludes \texttt{slistening} outright, so a live socket is never
unlinked.

\paragraph{Invariant 3 --- no two winners.}
At most one agent holds a non-absent claim on the path:
\[
  \# \{ a : AGENT | phase~a \in \{ bound, listening, serving \} \}
    \leq 1.
\]

\paragraph{Invariant 4 --- deadlock freedom.}
No reachable state leaves every agent unable to progress while some agent
is non-terminal.  The two-lock order (spawn-outer, bind-inner, never both
held by one agent, an agent holding one never requesting the other) makes
a cyclic wait impossible.

\paragraph{Invariant 5 --- lost-racer clean exit.}
An agent in \texttt{lostRace} can only reach \texttt{exited}, via
\texttt{Exit}, which changes nothing but its own phase.  By construction
\texttt{lostRace} is produced only by \texttt{BindFail} and
\texttt{LoseRaceLive} and consumed only by \texttt{Exit}; the \texttt{Exit}
schema frames \texttt{socketStatus}, \texttt{owner}, \texttt{spawnLock},
and \texttt{bindLock} unchanged.

% ============================================================
\section{Verification}
\label{sec:verify}
% ============================================================

Type-checking is a \texttt{fuzz} pass:
\begin{verbatim}
  fuzz -t docs/display_lifecycle.tex
\end{verbatim}

Invariants 1 and 3 are state properties, checked by asking ProB whether
their negation is \emph{reachable}.  A goal that is not found means the
invariant holds on every reachable state:
\begin{verbatim}
  # Invariant 1 (must report: goal NOT found)
  probcli docs/display_lifecycle.tex -model_check -nodead \
    -goal "card({a|a:AGENT & phase(a)=serving})>1" \
    -p DEFAULT_SETSIZE 2

  # Invariant 3 (must report: goal NOT found)
  probcli docs/display_lifecycle.tex -model_check -nodead \
    -goal "card({a|a:AGENT & phase(a):{bound,listening,serving}})>1" \
    -p DEFAULT_SETSIZE 2
\end{verbatim}

Invariant 4 is the deadlock check over the full reachable state space:
\begin{verbatim}
  # Invariant 4 (must report: no deadlock)
  probcli docs/display_lifecycle.tex -model_check \
    -p DEFAULT_SETSIZE 2
\end{verbatim}

Invariants 2 and 5 are transition/structural properties, established by
the shape of the operation schemas as argued in Section~\ref{sec:inv} and
corroborated by ProB's operation-coverage report (\texttt{CleanupStale}
never fires from a \texttt{slistening} state; \texttt{Exit} is the sole
consumer of \texttt{lostRace}).

% ============================================================
\section{Fidelity: reproducing the two-winner defect}
\label{sec:fidelity}
% ============================================================

A model that cannot reproduce the known defect proves nothing.  The
buggy variant is the current specification with one change: the bind-lock
serialisation is removed from the critical section, so cleanup and bind
no longer run under a single held lock.  Concretely, delete the guard
\texttt{a? $\in$ bindLock} from \texttt{CleanupStale} and \texttt{Bind}
(and, to let a second agent enter the window, from \texttt{Probe} and
\texttt{Listen}).  The two-winner interleaving then becomes reachable:

\begin{enumerate}
  \item Agent $x$ probes the absent path and binds it:
    \texttt{socketStatus = sbound}, \texttt{owner = \{x\}},
    \texttt{phase~x = bound}.
  \item Agent $y$ probes during the bind window and reads
    \texttt{sbound} as dead.
  \item Agent $y$ runs \texttt{CleanupStale}, unlinking $x$'s fresh
    socket: \texttt{socketStatus = sabsent}, \texttt{owner = $\emptyset$}
    --- while \texttt{phase~x} is still \texttt{bound}.
  \item Agent $y$ binds its own socket: \texttt{owner = \{y\}},
    \texttt{phase~y = bound}.
  \item Now \texttt{phase~x = bound} and \texttt{phase~y = bound}: two
    agents hold a non-absent claim, violating Invariant~3.
\end{enumerate}

Under the intact bind lock this interleaving is impossible: step~2
cannot occur because $y$ cannot acquire the bind lock while $x$ holds it,
and $x$ holds it from before step~1 until after it has listened and
released.  The reachability search for Invariant~3 therefore returns
\emph{goal not found} for the current design and \emph{goal found} --- the
trace above --- for the buggy variant.  That difference is the evidence
that the model detects the race the code was changed to prevent.

\end{document}
